% ============================================================================
% CAPÍTULO 1: MARCO TEÓRICO
% ============================================================================
% Este capítulo presenta la base teórica y conceptual que sustenta tu trabajo.
% Incluye conceptos clave, definiciones, y revisión de la literatura.
%
% TIPS PARA CITAS (formato APA, según la norma UTEM):
%
%   Cuando el autor forma parte de la narrativa:
%     Según \textcite{clave}, debido a la evolución...
%     Resultado: "Según Wilkinson (2001), debido a la evolución..."
%
%   Cuando el autor NO forma parte de la narrativa:
%     En un estudio reciente \parencite{clave}, se determinó que...
%     Resultado: "En un estudio reciente (Wilkinson, 2001), se determinó que..."
%
%   Cita textual corta (menos de 40 palabras):
%     Como señala el autor, ``texto citado'' \parencite{clave}.
%
%   Cita textual larga (más de 40 palabras):
%     \begin{quote}
%         Texto citado largo... \parencite{clave}
%     \end{quote}
% ============================================================================

\chapter{Marco Teórico}

% --- Introduce el capítulo brevemente ---
% En este capítulo se presentan los fundamentos teóricos que sustentan...

\section{Primer concepto o tema}

[Desarrolla aquí el primer concepto clave de tu marco teórico...]

% Ejemplo de cita:
% Según \textcite{apellido2020}, la definición de...

\subsection{Subtema 1.1}

[Desarrolla aquí un subtema...]

\subsection{Subtema 1.2}

[Desarrolla aquí otro subtema...]

\section{Segundo concepto o tema}

[Desarrolla aquí el segundo concepto...]

\section{Tercer concepto o tema}

[Desarrolla aquí el tercer concepto...]

% NOTA: Puedes agregar tantas secciones y subsecciones como necesites.
% Recuerda que cada sección debe tener un título claro y preciso.
